\documentclass[10pt,twocolumn]{article}
\usepackage[T1]{fontenc}
\usepackage[utf8]{inputenc}
\usepackage{lmodern}
\usepackage[margin=1.8cm,top=2.4cm,bottom=2cm,columnsep=0.6cm]{geometry}
\usepackage{fancyhdr}
\usepackage{titlesec}
\usepackage{booktabs}
\usepackage{amsmath,amssymb,amsfonts}
\usepackage{microtype}
\usepackage{xcolor}
\usepackage{mdframed}
\usepackage{parskip}
\usepackage{enumitem}
\usepackage{hyperref}

\definecolor{rulered}{RGB}{180,30,30}
\definecolor{boxgray}{RGB}{245,245,245}
\definecolor{keywordgray}{RGB}{100,100,100}

\pagestyle{fancy}
\fancyhf{}
\fancyhead[L]{\textsc{\small Claude Mag}}
\fancyhead[C]{\small\textit{Machine Learning Research Digest}}
\fancyhead[R]{\small 2026-09-07}
\fancyfoot[C]{\small\thepage}
\renewcommand{\headrulewidth}{0.8pt}

\titleformat{\section}{\normalsize\bfseries\color{rulered}}{}{0pt}{}%
  [\vspace{-4pt}\color{rulered}\rule{\columnwidth}{0.4pt}]
\titlespacing{\section}{0pt}{8pt}{4pt}

\hypersetup{colorlinks=true,urlcolor=rulered,linkcolor=black}

\begin{document}
\setlength{\parindent}{0pt}
\setlength{\parskip}{4pt}

{\small\color{keywordgray}\textbf{Keywords:}
  pixel-space diffusion \textbullet{}
  image generation \textbullet{}
  sparse attention \textbullet{}
  global refinement}\par\vspace{4pt}

{\LARGE\bfseries\raggedright
  Advanced Pixel Diffusion with Guided
  Sparse Global Refinement (PixSGR)}\par\vspace{4pt}

{\large\itshape\raggedright
  A Low-Channel Bottleneck Plus Sparse Cross-Patch
  Attention Brings Pixel Diffusion to Latent Diffusion
  Quality at Substantially Lower Cost}\par\vspace{6pt}

{\small\textbf{By} Weiyi You, Jinhua Zhang, Xingyu Zhou, Wei Long,
  Junyu Lou, Shuhang Gu}\par\vspace{2pt}

{\small\color{keywordgray}arXiv:2609.00798 \textbullet{} Preprint, September 2026}
\par\vspace{2pt}\noindent{\color{rulered}\rule{\columnwidth}{0.6pt}}\vspace{6pt}

Pixel-space diffusion offers higher fidelity and finer-grained control
than latent diffusion, but its computational cost has kept it behind
in practice. PixSGR resolves the core bottleneck with two innovations:
a supervised low-channel bottleneck that captures the low-dimensional
manifold of natural images before diffusion begins, and a sparse global
refinement step that restores structural coherence across patch boundaries
at sub-quadratic cost. On ImageNet $256\times256$, PixSGR matches latent
diffusion quality (FID 3.41 vs.\ 3.60 for LDM-4) while operating
entirely in pixel space.

\begin{mdframed}[backgroundcolor=boxgray,linecolor=rulered,linewidth=1pt,
    innerleftmargin=6pt,innerrightmargin=6pt,
    innertopmargin=6pt,innerbottommargin=6pt]
\textbf{\small AT A GLANCE}\par\vspace{4pt}
\begin{description}[leftmargin=0pt,labelsep=4pt,itemsep=2pt,parsep=0pt,
    font=\small\bfseries]
  \item[Problem:] Pixel-space diffusion is too computationally expensive
    and produces structural discontinuities at patch boundaries.
  \item[Why it matters:] Pixel-space models avoid latent encoder
    artifacts, enabling pixel-perfect fine-grained control.
  \item[Method:] Supervised low-channel bottleneck + patch-level
    denoising + sparse global cross-patch refinement.
  \item[Result:] FID 3.41 on ImageNet 256$\times$256, matching LDM-4
    with 40\% fewer FLOPs than dense pixel transformer baselines.
  \item[Evidence:] Ablations confirming sparse global refinement
    (FID 3.41 vs.\ 6.82 without it) and supervised bottleneck value.
\end{description}
\end{mdframed}

\section{The Big Picture}

Latent diffusion models (LDMs) dominate image generation by compressing
images into a learned latent space before applying the diffusion process.
This reduces computation dramatically but introduces a structural
limitation: the quality ceiling is bounded by the encoder-decoder and
artifacts from the latent representation propagate to the output.

Pixel-space diffusion avoids these artifacts by modelling the image
directly, but faces two challenges. First, a $256\times256$ RGB image
has 196{,}608 dimensions --- three orders of magnitude more than a typical
latent code. Second, existing pixel diffusion approaches either use
large patches (sacrificing fine detail) or process patches independently
(losing structural coherence across boundaries).

PixSGR addresses both challenges with a staged architecture that keeps
the effective dimensionality low while restoring global structure.

\section{Why Existing Approaches Fall Short}

\textbf{Dense pixel transformers} apply attention over all pixel tokens
jointly, achieving global coherence at the cost of $O(N^2)$ attention
where $N \approx 200$K --- computationally prohibitive at training scale.

\textbf{Patch-level diffusion (e.g.\ DiT in pixel space)} divides the
image into non-overlapping patches and denoises each independently.
This achieves linear cost but produces boundary artifacts: adjacent
patches generate independently without seeing each other's context,
leading to discontinuities in colour, texture, and large-scale structure.

\textbf{Hierarchical pixel diffusion} uses coarse-to-fine generation
but still confines refinement within each patch, missing long-range
interactions needed for global structural consistency.

\section{Core Mechanism}

PixSGR operates in three stages:

\textbf{Stage~1: Supervised Low-Channel Bottleneck.}
A supervised autoencoder compresses the image from $C=3$ channels to
$C'=4$ channels at the original spatial resolution. This is not a
latent encoder: spatial resolution is preserved. The bottleneck captures
the low-dimensional manifold of natural images in colour-space without
spatial downsampling.

\textbf{Stage~2: Patch-Level Diffusion.}
The diffusion process operates on non-overlapping patches of the
bottleneck representation. A patch-level transformer denoiser processes
patches in parallel --- linear in the number of patches.

\textbf{Stage~3: Sparse Global Refinement.}
After patch-level generation, a sparse set of refinement tokens is
positioned at patch boundaries and a random interior subset. Each
refinement token attends globally to all patch tokens via sparse
attention, restoring cross-patch coherence.

\section{Technical Formulation}

\subsection{Patch-Level Denoising Objective}

The denoising loss over patches is the standard DDPM objective:
\[
  \mathcal{L}_\mathrm{patch} = \mathbb{E}_{t, x_0, \varepsilon}
  \|\varepsilon - \varepsilon_\theta(x_t^\mathrm{patch}, t)\|^2
\]
where $x_t^\mathrm{patch}$ is the noisy patch at timestep $t$.

\subsection{Sparse Global Attention}

Let $P = \{p_1, \ldots, p_N\}$ be the full set of patch token positions
and $R = \{r_1, \ldots, r_M\}$ with $M \ll N$ be the sparse refinement
positions (straddling patch boundaries plus a random interior subset).
The refinement attention weight for refinement token $r_i$ over patch
token $p_j$ is:
\[
  A(r_i, p_j) =
  \frac{\exp\!\left(q(r_i)^\top k(p_j) / \sqrt{d}\right)}{Z_i}
  \;\;\text{for } j \in \mathcal{N}_\mathrm{local}(r_i) \cup \mathcal{N}_\mathrm{rand}(r_i)
\]
where $\mathcal{N}_\mathrm{local}$ is a fixed local neighbourhood and
$\mathcal{N}_\mathrm{rand}$ is a randomly sampled global subset. Total
attention cost is $O(N \log N)$ per refinement step.

\subsection{Guided Conditioning}

The refinement token hidden states are conditioned on the supervised
bottleneck representation at the corresponding spatial positions:
\[
  h_\mathrm{refine}(r_i)
  = \mathrm{Attention}(r_i, P)
  + \gamma \cdot \mathrm{Proj}\!\left(\mathrm{bottleneck}(r_i)\right)
\]
where $\gamma$ is a learnable scalar. This structural anchoring prevents
the refinement step from drifting from the coarse image structure.

\section{Learning or Inference Procedure}

Training proceeds in two phases:
\begin{enumerate}[itemsep=2pt,parsep=0pt]
  \item \textbf{Bottleneck pre-training.} Train the low-channel
    autoencoder with a reconstruction loss on ImageNet. The encoder
    produces the 4-channel spatial bottleneck; the decoder reconstructs
    the original image from patch-level outputs.
  \item \textbf{Joint diffusion training.} Train the patch-level
    denoiser and the sparse global refinement module jointly on the
    bottleneck representation, using the DDPM objective.
\end{enumerate}

At inference, the model runs the diffusion reverse process patch-by-patch,
then applies one sparse global refinement pass over the fully denoised
bottleneck before decoding to pixels.

\section{What the Guarantee Says}

The paper establishes that the sparse global attention achieves an
approximation error bound:
\[
  \|A_\mathrm{sparse} - A_\mathrm{dense}\|_F \leq
  \epsilon_\mathrm{local} + \epsilon_\mathrm{rand}
\]
where $\epsilon_\mathrm{local}$ decreases with the local neighbourhood
radius and $\epsilon_\mathrm{rand}$ decreases with the random sampling
budget. For the default configuration (16-token local neighbourhood,
5\% random sample), $\epsilon_\mathrm{rand} + \epsilon_\mathrm{local}$
is bounded under standard assumptions on the attention score distribution.

\section{Experimental Findings}

\textbf{Class-conditional generation, ImageNet $256\times256$:}

\begin{center}
\begin{tabular}{lcc}
\toprule
Model & FID$\downarrow$ & Params \\
\midrule
LDM-4 (latent diffusion) & 3.60 & 400M \\
DiT-XL/2 (latent)        & 2.27 & 675M \\
Pixel DiT (dense)        & 8.90 & 600M \\
\textbf{PixSGR (ours)}   & \textbf{3.41} & 550M \\
\bottomrule
\end{tabular}
\end{center}

PixSGR matches LDM-4 quality in pixel space with 8.5\% fewer parameters
than the dense pixel diffusion baseline, and reduces FLOPs per denoising
step by 40\%.

On FFHQ $256\times256$ (unconditional), PixSGR achieves FID 4.2
vs.\ 5.9 for the dense pixel baseline, with similar parameter counts.

\section{Ablations and Interpretation}

\textbf{Without sparse global refinement:} FID degrades from 3.41 to
6.82 on ImageNet --- the largest single ablation factor, confirming that
cross-patch coherence is the critical missing ingredient in prior pixel
diffusion.

\textbf{Without supervised bottleneck:} Convergence slows substantially
and final FID increases to 5.10; the bottleneck reduces the effective
dimensionality of the diffusion task from the start.

\textbf{Refinement token density:} FID improves as the sparse refinement
fraction increases from 5\% to 15\% of total patches, then plateaus.
The default of 10\% gives the best efficiency-quality trade-off.

\textbf{Guided vs.\ unguided refinement:} Adding the bottleneck
conditioning term (guided) reduces FID by 0.8 points over unguided
refinement, demonstrating that structural anchoring is beneficial.

\textbf{Local neighbourhood radius:} Increasing the local neighbourhood
from 4 to 16 tokens reduces FID by 0.4 points; beyond 16, gains are
marginal and cost increases.

\section*{Reference}

Weiyi You, Jinhua Zhang, Xingyu Zhou, Wei Long, Junyu Lou, and Shuhang Gu.
\textbf{Advanced Pixel Diffusion Model with Guided Sparse Global
Refinement.}
arXiv:2609.00798, September 2026.
\url{https://arxiv.org/abs/2609.00798}

\end{document}
